% Dieses Dokument ist eine Anleitung zur Benutzung von Circuitikz.
% Dabei werden einzelne Aspekte von Tikz und der Circuitikz Bibliothek
% Erklärt. Die gesamte Komposition eines Bildes wird hier nicht im Detail
% erläutert, dafür sind die jeweils einzelnen erstellten Bilder besser geeignet.

% Dieses Dokument kann direkt mit pdflatex kompiliert werden und
% ergibt ein PDF-Bild das mit Tikz gezeichnet wurde
% um einen Compile-Prozess zu starten, ist es sinnvoll, in der Kommandozeile
% 	latexmk -pdf anleitung_circuitikz.tex
% auszuführen um die Datei anleitung_circuitikz.pdf zu erhalten
% Die beim Kompilieren entstehenden Hilfsdateien können mit
% 	latexmk -C
% wieder entfernt werden.

% Als documentclass wird für gewöhnlich ein standalone Projekt verwendet.
% Dadurch ist keine "Blattgröße" vorgegeben und das Dokument wird nur so groß
% wie das darin gezeichnete Bild.
%
%\documentclass[border=2mm]{standalone}
%
% Leider lässt sich in einem standalone Projekt nur eine einzige Seite und somit
% auch nur ein einziges Bild erstellen.
% Für diese Anleitung, die mehrere Bilder beinhaltet, ist das etwas
% hinderlich. Deswegen wird % hier auf die article documentclass zurückegriffen.
\documentclass[a4paper,11pt]{article}

% Die Circuitikz Bibliothek bindet automatisch auch die Tikz Bibliothek ein.
\usepackage{circuitikz}

% Circuitikz unterstützt verschiedene Stile von Logikgattern.
% Hier wird der IEC Stil verwendet, der in Europa üblich ist.
% Alternativ könnte der amerikanische US Stil mit
% \usetikzlibrary{circuits.logic.US}
% verwendet werden.
\usetikzlibrary{circuits.logic.IEC}

% Um mathematische Ausdrücke im Bild verwenden zu können wird amsmath benötigt.
\usepackage{amssymb}

% Für die Darstellung von Einheiten wird siunitx verwendet.
\usepackage{siunitx}
% Hier wird die deutsche Lokalisierung verwendet, damit das Komma
% als Dezimaltrennzeichen verwendet wird.
\sisetup{locale = DE}

% Beginn des Dokumentes, in dem sich das Bild befindet.
\begin{document}
% Ab hier werden vereinzelt Tikz Bilder mit Circuitikz gezeichnet und auf den Seiten verteilt
% Die Bilder sind alle ausführlich dokumentiert, damit der Erstellungsprozess nachvollziehbar ist.

%%%%%%%%%%%%%%%%%%%%%%%%%
% Kapitel 1: Grundlagen %
%%%%%%%%%%%%%%%%%%%%%%%%%

% Jedes Tikz Bild beginnt mit einer 'tikzpicture' Umgebung.
% Eine Circuitikz Umgebung ist eine spezielle Tikz Umgebung, in der
% die Circuitikz Bibliothek automatisch eingebunden wird.
% Hierbei wird mit dem [european] flag vorgegeben, dass der europäische Stil verwendet wird.
\begin{circuitikz}[european]
	% Eine Tikz Umgebung besteht aus einer Reihe von '\draw' und '\node' Anweisungen.
	% wobei auch nodes in einer Draw Anweisung erstellt werden können. Tikz ermöglicht es,
	% das Ziel auf sehr viele verschiedene Weisen zu erreichen.
	% Deswegen beschreibt dieses Dokument nur einen möglichen Weg,
	% versucht dabei aber, seinen eigenen Stil konsequent durchzuhalten.

	% Tikz unterscheidet zwischen folgenden Elementen:
	% 1. Koordinaten: Jedes Tikz-Bild existiert in einem fiktiven Koordinatensystem.
	%    Koordinaten werden in runden Klammern angegeben, z.B. (0,0) für den Ursprung.
	%    Koordinaten können einen willkürlichen Bereich aufspannen, der sowohl positive als auch
	% 	 negative Werte enthalten kann. Dabei muss der Ursprung auch nicht die Mitte des Bildes sein
	%    sondern kann sich an einer willkürlich gewählten Stelle befinden.
	%	   Alle Koordinaten in Tikz sind in gewisser Weise willkürlich, da es kein fest
	%		 vorgeschriebenes Konzept dafür gibt, wie die Koordinaten gewählt werden sollen.
	%    Es ist aber sinnvoll, sich an einem Raster zu orientieren, um die Übersicht zu behalten.
	% 2. Nodes: Nodes sind Objekte, die an einer bestimmten Koordinate positioniert werden.
	%		 Es gibt verschiedene Arten von Nodes, z.B. Textnodes, die einfach nur Text an einer
	%		 bestimmten Koordinate darstellen. Es gibt aber auch Nodes, die komplexere Objekte
	%		 darstellen, z.B. Logikgatter, Widerstände, Kondensatoren, usw. Diese Knoten werden z.T.
	%    durch die Circuitikz Bibliothek bereitgestellt.
	%		 Ein Knoten kann auch einen Namen bekommen, über den er später referenziert werden kann.
	%		 Dieser Name ist wie eine Variable, die die Koordinaten des Knotens speichert.
	% 	 Für gewöhnlich hat ein Knoten fest definierte Ankerpunkte, die ebenfalls
	%    referenziert werden können.
	% 3. Pfade: Pfade sind Linien, die zwei oder mehr Koordinaten oder Knoten miteinander verbinden.
	% 	 Dabei können Pfade auch verschiedene Stile und Formen haben. Manche Objekte in der Circuitikz
	%    Bibliothek werden auch durch Pfade realisiert, z.B. Widerstände, Kondensatoren, Dioden.

	% Zu Beginn schauen wir uns mal ein einfaches Beispiel an: Wir zeichnen eine Linie zwischen zwei Koordinaten.
	\draw
		(0,0) -- (4,0)
	;
	% Die \draw Anweisung beginnt mit dem \draw Befehl und endet mit einem Semikolon. Dazwischen stehen
	% die Anweisungen, die ausgeführt werden sollen. In diesem Fall wird eine Linie zwischen
	% den Koordinaten (0,0) und (4,0) gezeichnet. Die Linie wird durch zwei Bindestriche '--' erzeugt.
	% Ein Koordinaten-Paar besteht aus einer X- und einer Y-Koordinate, die durch ein Komma getrennt sind.
	% Dabei verläuft die X-Achse horizontal und mit zunehmendem Wert nach rechts.
	% Die Y-Achse verläuft vertikal und mit zunehmendem Wert nach oben.
	% Die hier gezeichnete Linie verläuft also parallel zur X-Achse.

\end{circuitikz} \\
Bild 1: Eine einfache Linie zwischen zwei Koordinaten.

%%%%%%%%%%%%%%%%%%%%%
% Kapitel 2: Linien %
%%%%%%%%%%%%%%%%%%%%%

% Das nächste Bild soll die Koordinaten eines Tikz Bildes verdeutlichen.
\begin{circuitikz}[european]
	% Wie bereits gesagt können die Koordinaten eines Tikz Bildes willkürlich gewählt werden.
	% Das Bild würde identisch aussehen, wenn stattdessen folgende Anweisungen verwendet würden:
	\draw
		(-2,1) -- (2,1)
	;
	% Die Linie bleibt gleich lang und das Bild erscheint identisch.
	% Im PDF-Dokument sollten also zwei Linien zu sehen sein, die exakt übereinander liegen.

	% An dieser Stelle ist noch ein Hinweis über die Linie sinnvoll, die mit '--' erzeugt wird.
	% Dabei handelt es sich um einen sogenannten Pfad.
	% Mit Pfaden werden zwei Koordinaten verbunden.
\end{circuitikz} \\
Bild 2: Dieselbe Linie an der selben Position, obwohl die Koordinaten anders gewählt wurden.

\begin{circuitikz}[european]
	% Wir können die Art der Verbindungslinie anpassen. Für gewöhnlich wird eine gerade Linie die
	% beeiden Koordinaten direkt verbinden. Das kann zur Folge haben, dass die Linie auch diagonal verläuft.
	\draw
		(0,0) -- (4,2)
	;
	% In diesem Fall verläuft die Linie diagonal von (0,0) nach (4,2).
	% In manchen Fällen ist das aber nicht gewünscht. Stattdessen kann mit '-|' und '|-'
	% angegeben werden, dass die Linie zuerst horizontal und dann vertikal verlaufen soll bzw. umgekehrt.
	% Das Ergebnis sieht dann so aus:
	\draw
		(5,0) |- (9,2) % Zuerst vertikal dann horizontal
	;
	\draw
		(10,2) |- (12,0) % Zuerst horizontal dann vertikal
	;
	% Man sieht hier übrigens auch, dass sehr viele draw Befehle in einem circuitikz Bild stehen können.
\end{circuitikz} \\
Bild 3: Diagonale, Horizontal-Vertikale Linien und Vertikal-Horizontale Linien.

\begin{circuitikz}[european]
	% Koordinaten können auch relativ zueinander angegeben werden. Dazu wird die Koordinate
	% in runde Klammern gesetzt und mit einem '++' vorangestellt.
	% Das bedeutet, dass die Koordinate in Klammern zur vorherigen Koordinate addiert wird um
	% die Ziel-Koordinate zu erhalten.
	% Dieses Beispiel hier zeigt ein paar Linien, die mit relativen Koordinaten zueinander gezeichnet wurden.
	\draw
		(0,0) -- ++ (4,2) -- ++ (2,-1) -- ++ (-1,-3) -- (0,0)
	;
	% Die Linie beginnt bei (0,0) und verläuft dann 4 Einheiten nach rechts und 2 Einheiten nach oben.
	% Von dort aus verläuft sie 2 Einheiten nach rechts und 1 Einheit nach unten.
	% Dann verläuft sie 1 Einheit nach links und 3 Einheiten nach unten.
	% Zum Schluss verläuft sie wieder zurück zum Ursprung (0,0). Hier wurde keine relative Koordinate verwendet.
	% In Tikz kann also jederzeit zwischen absoluten und relativen Koordinaten gewechselt werden.
	% Das Ergebnis ist eine geschlossene Linie, die wieder am Ausgangspunkt endet.
	%
	% Ein weiterer Interessanter Aspekt ist, dass die Linien in einer einzigen draw Anweisung stehen. 
	% Tikz überlässt dem Benutzer die Wahl, ob er viele einzelne draw Anweisungen verwenden möchte
	% oder alles in eine einzige draw Anweisung packt.
	% Die Relativen Kordinaten können ebenfalls mit '-|' und '|-' kombiniert werden.
	\draw
		(6,0) -- ++ (2,2) |- ++ (2,-1) -| ++ (2,1) -- ++ (-6,-4)
	;
\end{circuitikz} \\
Bild 4: Linien, die mit relativen Koordinaten gezeichnet wurden.

% Damit die Arbeit mit Koordinaten nicht zu verwirrend wird,
% gibt es eine einfache Möglichkeit, sich während der Arbeit am Bild
% ein Koordinatensystem zeichnen zu lassen.
\begin{circuitikz}[european]
	% Dazu kann der \draw Befehl mit Parametern aufgefordert werden,
	% das Koordinatensystem im Stil von 'help lines'
	% zu zeichnen:
	\draw[help lines] (0,0) grid (3,2);
	% Dieser Befehl zeichnet ein Grid um das Koordinatensystem darzustellen
	% Nachdem das Bild fertig gezeichnet wurde, kann man diese Zeile wieder auskommentieren
	\draw (0.5,0) -- ++ (2,1.5);
	% Koordinaten müssen keine ganzen Zahlen sein. Dabei wird der Punkt '.'
	% als Dezimaltrennzeichen verwendet.
\end{circuitikz} \\
Bild 5: Ein Koordinatensystem, das die Arbeit erleichtert


%%%%%%%%%%%%%%%%%%%%%
% Kapitel 3: Knoten %
%%%%%%%%%%%%%%%%%%%%%

% Bisher haben wir uns Koordinaten und Pfade angeschaut.
% Pfade verbinden **zwei** Koordinaten.
% Knoten sind da anders: Sie befinden sich an **einer** Koordinate.

% Als Knoten lässt sich jedes Objekt bezeichnen, das "komplizierter als eine Linie" ist.
% Schauen wir uns dazu mal einfache Knoten an:
\begin{circuitikz}[european]
	% Knoten können einerseits über den \node Befehl erstelle werden.
	% Dabei gibt der 'at' Operator die Position des Knoten an:
	\node at (0,0) {A simple text node};
	% Wir können dem node-Befehl ein paar parameter übergeben.
	% Beispielsweise können wir für den Knoten einen Rand zeichnen lassen
	\node[draw] at (0,-1) {A text node with a border};
	% Die angegebene Koordinate bezieht sich dabei jeweils auf den Mittelpunkt des Knoten.
	% Wir können die Knoten aber auch an einem anderen Ankerpunkt als dem Mittelpunkt ausrichten
	% uns damit dafür sorgen, dass der Text sich z.B. links oder rechts von einer vorgegebenen Koordinate
	% befindet
	\node[left, draw] at (4.5,-2) {A left node};
	\node[right, draw] at (5.5,-2) {A right node};
	% Drunter und drüber sind auch zwei Möglichkeiten
	\node[above] at (5,-3) {A node above the specified coordinate};
	\node[below, draw] at (5,-3) {A node below the specified coordinate};
	% zur besseren Darstellung wurde über den 'draw' Parameter der Rand des Knoten gezeichnet,
	% damit die Text-Knoten im gerenderten Bild besser auseinander gehalten werden können.
\end{circuitikz} \\
Bild 6: Einfache Text-Knoten

% Auf Dauer kann es etwas anstrengend werden, wenn jeder Knoten in einem eigenen
% \node Befehl platziert werden kann. Wie bereits erwähnt gibt es mit Tikz mehrere Wege
% zum Ziel. Knoten können auch genau so gut in einem \draw Befehl platziert werden:
\begin{circuitikz}[european]
	\draw
		(0,0) node[draw] {Noch ein Text-Knoten}
	;
	% Fürs erste wirkt da vielleicht etwas komisch, da es plötzlich zwei
	% verschiedene Möglichkeiten gibt, Knoten zu erstellen, aber
	% der \draw Befehl wird sich als die weitaus flexiblere Möglichkeit
	% erweisen, Knoten zu platzieren.
	% Das liegt daran, dass im \draw Befehl relative Koordinaten und
	% Pfade verwendet werden können und so ein Knoten abhängig von der Position eines
	% anderen Knoten platziert werden kann.
	% Mit dem \node Befehl wäre das etwas aufwendiger:
	\draw
		(1,-1) node[draw, right] {Rechter Text-Knoten}
		% jetzt können wir einen relativen Sprung von der vorherigen Koordinaten aus vornehmen
		++ (0,-0.5) node[below] {Und noch ein Knoten darunter}
		% Dabei ist folgendes zu beachten:
		% Die angegebene Koordinate (1,-1) wurde auf den linken Ankerpunkt des ersten
		% knoten gelegt. Der Text befindet sich rechts von der Koordinate
		% Danach wurde ein Sprung um 0.5 Einheiten nach unten vorgenommen um einen weiteren
		% knoten an dieser Stelle zu platzieren. Für diesen Knoten wurde aber der oberste Punkt
		% als Anker angegebene (below), was dazu führt, dass der Text auf der horizontal-Achse
		% um diese Koordinate herum seinen Mittelpunkt hat. Deswegen sind die Texte nicht nur vertikal,
		% sondern auch horizontal verschoben.
	;
\end{circuitikz} \\
Bild 7: Noch mehr Knoten

% In dieser Anleitung wird für die meisten Bilder der \draw Befehl verwendet um Knoten
% zu zeichnen. Dies liegt darin, dass der \draw Befehl mehr Flexibilität in der Platzierung
% der Knoten ermöglicht. Weiterhin ist es damit auch möglich, das gesamte Bild in einem
% einzigen (vielleicht etwas langem) \draw Befehl auszudrücken.

%%%%%%%%%%%%%%%%%%%%
% Kapitel 4: Anker %
%%%%%%%%%%%%%%%%%%%%

% Vorhin haben wir bereits bei den Text-Knoten Ankerpunkte verwendet, um den Text über, unter,
% links und rechts von einer Koordinate zu platzieren.
% Diese 4 Schlüsselwärter sind für Text-Knoten in Tikz gängig, aber es gibt für alle Knoten
% noch allgemeinere Ankerpunkte, die etwas mehr Flexibilität bereitstellen.
%
% Die Ankerpunkte eines Knoten sind wie auf einem Kompass verteilt:
%
%		north west			north				north east
%				*-------------*-------------*
%				|														|
%				|														|
%				|														|
%				|														|
%  west *							* center			* east
%				|														|
%				|														|
%				|														|
%				|														|
%				*-------------*-------------*
%		south west			south				south east
%
% Diese Ankerpunkte sind im eigentlichen Sinne Koordinaten, die referenziert werden können.
% Das ganze soll an ein paar Beispielen erklärt werden:
\begin{circuitikz}[european]
	% Ab hier werden die Bilder alle in einem einzigen Draw Befehl gezeichnet,
	% der zum Schluss mit einem einzigen Semikolon abschließt.
	\draw
		(0,0) -- (1,0)  % Zu erst eine Linie
		node[anchor=west] {Linie}  % und ein Text-Knoten rechts davon

		% Dabei ist die zweite Zeile identisch zu der Anweisung
		% node[right] {Linie}
		(0,-1) -- (1,-1)  % Zu erst eine Linie
		node[right] {Linie}

		% Bis jetzt war alles noch schön und gut. Jetzt beginnt der eigentliche Spaß mit Knoten:
		% Wie bereits erwähnt, sind die Ankerpunkte Koordinaten. Diese Koordinaten können auch
		% referenziert werden um z.B. weitere Knoten an ihnen zu platzieren oder Linien zu diesen
		% Ankerpunkten zu ziehen.
		% Damit wir das machen können, müssen wir mit der vollständigen Syntax eines Knoten
		% vertraut werden:
		%
		%		node[parameter] (Referenz) {Text}
		% 
		% - Die Parameter bestimmen die Knoten-Art und geben damit vor, wie der Knoten gezeichnet
		%   werden soll. Dazu hatten wir bereits die Beispiele left, right, ... und die Option
		%	  anchor=... verwendet.
		% - Über die Referenz wurde bisher noch nichts erwähnt. Die Referenz ist eine Möglichkeit,
		%		den Knoten in einer **Variablen** zu speichern und später über diese Variable
		%		wieder zu referenzieren. Dazu folgen gleich ein paar Beispiele
		% - Der Inhalt war bei Text-Knoten bisher immer der darzustellende Text
		%
		% Schauen wir uns mal ein par Beispiele dazu an, wie mit Referenzen gearbeitet wird:
		% (Hinweis: Wir befinden uns immer noch in der \draw Anweisung)
		(0,-2) node[right] (R1) {Knoten R1}
		% Hiermit haben wir einen Knoten angelegt, den wir über die Variable (R1) ansprechen können
		% jetzt wollen wir rechts davon eine Linie zeichnen
		(R1.east) -- ++ (1,0)
		% und dann einen weiteren Knoten platzieren
		node[anchor=north west] (NW1) {Knoten NW1}
		% Wenn wir wollen, können wir jetzt eine Linie quer durch den Knoten NW1 ziehen
		(NW1.south west) -- (NW1.north east)
		% Man sieht also: Die Ankerpunkte von Knoten sind wie Variablen und folgen auch derselben
		% Syntax: Sie werden in runden Klammern (...) referenziert.
		% Zum Schluss können wir auch den Rahmen um einen Knoten selber zeichnen:
		(-1,-3) node[anchor=center] (C) {Rahmen, aber ohne den 'draw' Parameter}
		(C.north west) -| (C.south east) -| (C.north west)
		% Hierbei handelt es sich wieder um einen geschlossenen Streckenzug wie er bei den
		% ersten Beispielen mit Linien gezeichnet wurde.
	;
\end{circuitikz} \\
Bild 8: Die Verwendung von Ankern

%%%%%%%%%%%%%%%%%%%%%%%%%%%%%%%%
% Kapitel 5: Circuitikz Knoten %
%%%%%%%%%%%%%%%%%%%%%%%%%%%%%%%%

% Ab jetzt schauen wir uns einmal an, welche Knoten die Circuitikz Bibliothek mitliefert.
%
% Circuitikz stellt eine große Anzahl an Knoten-Objekten bereit um Logikgatter, Transistoren und
% andere elektronische Bauteile zu zeichnen.
% Dabei sind alle Bauteile ausführlich in der Circuitikz Dokumentation beschrieben:
%		https://ctan.dcc.uchile.cl/graphics/pgf/contrib/circuitikz/doc/circuitikzmanual.pdf
% Hierbei ist es sinnvoll, sich nicht die Minimalbeispiele in der Dokumentation anzuschauen,
% sondern im Kapitel "The Components: List" nach den Bauteilen zu suchen die platziert werden sollen.
% Dort werden auch alle Ankerpunkte für die Komponenten definiert.
% Die Ankerpunkte sind besonders wichtig, da sie die Stellen sind, an denen die Bauteile
% miteinander verbunden werden.
%
% Schauen wir uns dazu erst einmal ein paar Beispiele an:
\begin{circuitikz}[european]
	\draw
		% zuerst platzieren wir einen NPN Bipolar-Transistor:
		(0,0) node[npn, anchor=B] (T1) {T1}
		% dieser Bekommt die Bauteilbeschriftung T1, die in den geschweiften Klammern
		% {...} angegeben wird.
		% Der Ankerpunkt 'B' (Basis) wird auf die Koordinate (0,0) gelegt.
		% Die anderen Ankerpunkte sind der Collektor (C) und der Emitter (E).
		% Über die Variable kann der Transistor (T1) später referenziert werden.
		%
		% jetzt zeichnen wir eine Linie vom Collektor (C) des Transistors zur Basis eines
		% weiteren Transistors:
		(T1.C) -- ++ (2,0) node[npn, anchor=B] (T2) {T2}
		%
		% Wenn in elektrischen Schaltplänen Linien, die sich kreuzen, miteinander verbunden sind,
		% wird dies häufig durch einen Punkt dargestellt. Dieser kann durch einen 'circ' node erstellt werden:
		(0,-2) -| (3,-4)
		(3,-2) node[circ] {} -- (5,-2)
		% Der 'circ' node hat keine Bauteilbeschriftung, die Geschweiften Klammern {...} sind leer.
		% Dennoch ist es verpflichtend, die Geschweiften Klammern nach der Platzierung eines Nodes
		% zu setzen. Die runden Klammern sind optional, die geschweiften sind verpflichtend.
		% Andernfalls resultiert dies in einer Fehlermeldung beim Kompilieren.
		%
		% Es gibt den Punkt auch in der 'ocirc' Version, bei der das Innere weiß ist.
		% Diese Version is sinnvoll, wenn es darum geht, bei einer logischen Schaltung einen
		% Eingang oder einen Ausgang zu kennzeichnen.
		% Alternativ kann man damit bei einem Logikgatter auch eine Negation des Ausgangssignals
		% kennzeichnen. Das vorherige Beispiel wird nun um einen Eingangs-Port bzw. einen Ausgangs-Port
		% mit der Beschriftung A, B und Y erweitert (die Ports sind nur fiktiv und haben in dem
		% Beispiel keine Bedeutung)
		(0,-2) node[ocirc, anchor=west] {} node[left] {A}
		(3,-4) node[ocirc, anchor=south] {} node[below] {B}
		(5,-2) node[ocirc, anchor=east] {} node[right] {Y}
		% Warum wurde in diesem Beispiel für jeden 'ocirc' Node der Ankerpunkte genau so gewählt?
		% Schauen wir uns dazu mal die erste Zeile an: Die Linie beginnt bei (0,-2) und geht nach
		% rechts. Wenn wir uns nun den 'ocorc' Node weg denken, dann bleibt nur noch die Beschriftung
		% {A} rechts von der beginnenden Linie übrig. Das stimmt schon mal. Falls wir jetzt aber
		% den Kreis ebenfalls rechts von der beginnenden Linie setzen, haben wir das Problem,
		% dass der 'ocirc' node und der Buchstabe "A" sich sehr nahe kommen bzw. überlappen können.
		% Das würde nicht schön aussehen. Deshalb wird der 'ocirc' node **rechts** von der beginnenden
		% Linie platziert. Weil das Innere des Kreises weiß aufgefüllt wird (und nicht transparent
		% ist), wird damit der Beginn der Linie übermalt. Das Bild sieht damit am besten aus.
	;
\end{circuitikz} \\
Bild 9: Eine Sammlung an einfachen Circuitikz Bauteilen

% Nun werfen wir den Blick mal auf eine bestimmte Gruppe an Circuitikz Bauteilen: Logikgatter.
% Die Symbole für Logikgatter sind historisch gewachsen, weshalb es mehrere, zum Teil
% konkurrierende Standards dafür gibt. Leider ist Circuitikz davon nicht
% verschont geblieben.
% Für die europäischen Symbole gibt es in Kapitel 4.22.3 der Circuitikz Dokumentation eine
% genauere Erklärung zu den Bauteilen.
\begin{circuitikz}[european]
	\draw
		% Die '<gatter> port' Node Elemente sind eine ältere Gruppe an europäischen Symbolen.
		% Diese Bauteile sind eine sehr alte Gruppe, die aus Kompatibilitätsgründen für immer
		% ein Teil von Circuitikz bleiben wird. Eigentlich sollte sie durch einen Nachfolger
		% abgelöst werden, der im Verlauf dieser Anleitung genauer beschrieben wird, aber diese
		% Symbole sind für viele Zeichnungen immer noch besser geeignet als die neue Version.
		% Aus diesem Grund wurde für die Praktikums-Anleitung maßgeblich auf die 'port style
		% logic gates zurück gegriffen.
		(0,0) node[and port, anchor=out] {}
		% Ein 'port style logic gate' (so heißen diese Node-Typen) hat immer folgende Ankerpunkte:
		% - 'in 1' und 'in 2' sind die Eingänge auf der Linken Seite des Gatters
		% - 'out' ist der Ausgang auf der rechten Seite des Gatters
		% zusätzlich gibt es noch Ankerpunkte direkt am Gatter an den jeweiligen Ein- und ausgängen:
		% - 'bin 1' und 'bin 2' sind die Ankerpunkte für die Eingänge
		% - 'bout' ist der Ausgang des Gatters
		% die mit 'b' beginnenden Ankerpunkte werden demnächst wichtig, wenn es darum geht,
		% die Ausgänge zu negieren:
		%
		% Die Tatsache, dass diese Logikgatter veraltet sind, bringt ein paar Nachteile mit sich.
		% Einer dieser Nachteil wird sichtbar, wenn man sich z.B. das NAND Gatter anschaut:
		(3,0) node[nand port, anchor=out] {}
		% Die Grafische Darstellung zeigt den negierten Ausgang nicht mit einem Punkt an, wie es
		% eigentlich in der IEC norm vorgesehen ist. Stattdessen befindet sich ein schräger Strich am
		% Ausgang des Gatters. Diese Darstellung wird eigentlich nicht mehr verwendet.
		% Für die Praktikums-Anleitung wurde deshalb auf einen Workaround zugegriffen:
		% Die Logikgatter wurden in der nicht-negierten Form gezeichnet und der Kreis für die 
		% Negation wurde nachträglich mit einem 'ocirc' node hinzugefügt. Für das NAND würde das
		% folgendermaßen aussehen:
		(6,0) node[and port, anchor=out] (NAND) {}
		(NAND.bout) node[ocirc, right] {}
		% Dieser Trick funktioniert für alle logikgatter. Wenn man ein NOT Gatter zeichnen möchte,
		% muss man dafür das 'buffer port' Gatter negieren:
		(9,0) node[buffer port, anchor=out] (NOT) {}
		(NOT.bout) node[ocirc, right] {}
		% Warum wird für die Platzierung des NOT Gatters als Anker eigentlich 'out' verwendet?
		% Der Grund ist, dass in diesem Beispiel die Gatter in einer Reihe gezeichnet
		% werden sollen. Das NOT Gatter hat aber nur einen Eingang. Wenn wir den Ankerpunkt
		% 'in 1' verwenden würden, wäre ein Höhen-Abstand zwischen dem NOT Gatter und den
		% vorherigen Gattern sichtbar, weil der eine Eingang auf einer anderen Höhe liegt.
		% Deswegen wird der Ankerpunkt 'out' verwendet, der für alle Gatter auf derselben Höhe liegt.
	;
\end{circuitikz} \\
Bild 10: Ein paar Logikgatter

% Wie bereits erwähnt gibt es eine zweite Gruppe and Logikelementen, die für das Praktikum
% aber selten verwendet wurde. Dies wird hier kurzgefasst erklärt.
\begin{circuitikz}[european, circuit logic IEC]
	\draw
		% Die '<gatter> gate' Node Elemente sind eine alternative Gruppe an Logikgattern.
		% Damit die Zeichnung auch erfolgreich kompiliert, muss der 'circuit logic IEC' Parameter
		% in der circuitikz Umgebung angegeben werden! (siehe vier Zeilen oben).
		(0,0) node[and gate, anchor=output] {}
		% Die Anker heißen 'input 1', 'input 2' und 'output', d.h. der Name wird ausgeschrieben.
		% Diese Gruppe stellt die Negierung eines Ausgangs richtig dar:
		(3,0) node[nand gate, anchor=output] (NAND) {}
		(6,0) node[not gate, anchor=output] (NOT) {}
		% Ein Nachteil an diesen Gates ist aber, dass keine Linien an den Ein- und Ausgängen
		% gezeichnet werden. Dazu müssen von den Ankerpunkte aus manuell Linien gezeichnet werden.
		% Hier einmal für das NOT Gatter:
		(NOT.output) -- ++ (0.25,0)
		(NOT.input) -- ++ (-0.25,0)
		% und hier für das NAND Gatter:
		(NAND.input 1) -- ++ (-0.25,0)
		(NAND.input 2) -- ++ (-0.25,0)
		(NAND.output) -- ++ (0.25,0)
	;
\end{circuitikz} \\
Bild 11: Ein paar anders aussehende Logikgatter

%%%%%%%%%%%%%%%%%%%%%%%%%%%%%%%%%%%%%%%%%
% Kapitel 5: komplexe Logik-Schaltungen %
%%%%%%%%%%%%%%%%%%%%%%%%%%%%%%%%%%%%%%%%%

% Meistens bestehen logische Schaltungen aus mehreren Gattern, die miteinander verbunden sind
% Dazu soll hier einmal ein Beispiel gezeigt werden, um zu erklären, wie man solche Schaltungen
% sinnvoll mit Circuitikz zusammenstellen kann.
%
% Für gewöhnlich wird eine Schaltung von links nach rechts "gelesen". Die Eingänge befinden sich
% links und die Ausgänge rechts. Die meisten Schaltungen, die im Praktikum verwendet werden,
% haben viele Eingänge auf der linken und wenige Ausgänge (meistens einen) auf der rechten Seite.
% Das bedeutet automatisch, dass die Schaltungen von rechts nach links eine Baum-Struktur
% aufspannen. Es ist bei manchen Schaltungen also sinnvoll, beim Zeichnen von rechts nach links
% vorzugehen, damit die Baum-Struktur verwendet werden kann um die Bauteile nahezu automatisch
% anzuordnen.
% Dies ist natürlich nicht immer der Fall, aber bei diesem Beispiel wird dieses Phänomen ausgenutzt.
%
% Dazu zeichnen wir einmal den "2-Bit Komparator" aus Aufgabe 2.3 des Praktikumsversuchs.
\begin{circuitikz}[european, circuit logic IEC]
	\draw
		% wir beginnen mit dem AND Gatter ganz rechts. Die Wahl des Ursprungs ist bei Tikz egal.
		(0,0) node[and port, anchor=out] (AND) {}
		% Den Ausgang des AND Gatters beschriften wir mit einem Y und zeichnen einen Kreis.
		% Wie bereits vorher erwähnt wird für den Kreis ein ungewöhnlicher Anker verwendet,
		% damit er nicht zu nah am Buchstaben Y sitzt.
		(AND.out) node[ocirc, left] {} node[right] {Y}

		% Im nächsten Schritt zeichnen wir die beiden Eingänge des AND Gatters. Jeder Eingang ist mit
		% einem XNOR verbunden (einem negierten XOR). Dazu verwenden wir jeweils ein XOR Gatter und
		% malen den Punkt zur Negierung des Ausgangs manuell dazu.
		% Fangen wir mal mit dem Oberen Gatter an:
		(AND.in 1) -| ++ (0,0.5) node[xor port, anchor=out] (XNOR1) {}
		% Der Ausgang wird negiert:
		(XNOR1.bout) node[ocirc, right] {}
		% Diese beiden Eingänge werden nun beschriftet, ähnlich wie der Ausgang:
		(XNOR1.in 1) node[ocirc, right] {} node[left] {A1}
		(XNOR1.in 2) node[ocirc, right] {} node[left] {B1}
		
		% Dasselbe nochmal: Zuerst das XNOR-Gatter:
		(AND.in 2) -| ++ (0,-0.5) node[xor port, anchor=out] (XNOR2) {}
		(XNOR2.bout) node[ocirc, right] {}
		% Beschriftung der Ausgänge:
		(XNOR2.in 1) node[ocirc, right] {} node[left] {A2}
		(XNOR2.in 2) node[ocirc, right] {} node[left] {B2}
		% Ob die Ein- und Ausgänge mit einem 'circ' oder 'ocirc' Node versehen werden,
		% ist Geschmackssache. Manchmal sie die Schaltung damit besser aus und manchmal verwirrt das.
	;
\end{circuitikz} \\
Bild 12: Die erste komplizierte Schaltung

%%%%%%%%%%%%%%%%%%%%%%%%%%%%
% Kapitel 6: Pfad-Elemente %
%%%%%%%%%%%%%%%%%%%%%%%%%%%%

% Nicht alle Circuitikz-Bauteile sind Node-Elemente. Wenn ein Bauteil nur zwei Anschlüsse hat, dann
% ist eine Beschreibung im Stil von "Zeichne einen Widerstand von Koordinate A nach Koordinate B"
% kompakter als die Platzierung eines node-Widerstandes, welcher dann mit Linien verbunden
% werden muss.
% Solche Bauteile werden als Pfad-Elemente bezeichnet. In der Praktikums-Anleitung betrifft das
% nur Widerstände und Dioden.
\begin{circuitikz}[european]
	\draw
		% Ein Widerstand ist ein Pfad-Element, welches von einer Koordinate zu einer anderen
	 	% Koordinate gezeichnet wird. Dabei kann der Widerstand mit einer Bauteilbeschriftung
		% versehen werden.
		(0,0) to[resistor, l=$R_1$] ++ (2,0)
		% Anstelle der Bauteilbeschriftung kann das Bauteil auch mit einem Wert versehen werden
		to[resistor, l=$\SI{10}{\ohm}$] ++ (2,0)
		% Weitere Informationen über Widerstände sind im Circuitikz Manual in Kapitel 4.2 zu finden.

		% Dioden verhalten sich ähnlich. Die Polarisation der Diode wird durch die Pfadrichtung
		% (Quelle- -> Ziel-Koordinate) angegeben.
		to[diode, l=D1] ++ (2,0)
		% Dioden lassen sich auch einfärben, insbesondere LEDs:
		to[led, fill=red, -] ++ (1.5,0)
		
		% ground und Vcc sind **keine** Pfad-Elemente:
		node[ground] {}
		(0,0) node[vcc] {\SI{5}{\volt}}

		% Wenn sich zwei Pfade kreuzen und an dieser Stele elektrisch miteinander verbunden sein sollen,
		% wird dies in der Schaltung mit einem Punkt dargestellt. Damit diese Punkte nicht manuell
		% als 'circ' Node hinzugefügt werden müssen, gibt es die Möglichkeit, bei Pfaden automatisch
		% Punkte an der Start- oder End-Koordinate zu zeichnen:
		(2,0) to[resistor, l=$R_2$, *-] ++ (0,-2)
		node[ground] {}
		% - Der '*-' Parameter sorgt dafür, dass am Startpunkt des Pfades ein Punkt gezeichnet wird.
		% - Der '-*' Parameter sorgt dafür, dass am Endpunkt des Pfades ein Punkt gezeichnet wird.
		% - Mit '*-*' werden an beiden Enden Punkte gezeichnet.
	;
\end{circuitikz} \\
Bild 13: Widerstände und Dioden in Circuitikz

%%%%%%%%%%%%%%%%%%%%%%%%%%%%%%%%%%
% Kapitel 7: Schnitt-Koordinaten %
%%%%%%%%%%%%%%%%%%%%%%%%%%%%%%%%%%

% Ein sehr sinnvolles Feature in Tikz sind Schnittpunkt-Koordinaten.
% Dazu schauen wir uns einmal folgendes Problem an:
\begin{circuitikz}[european]
	\draw
		% Wir wollen den 1-Bit Multiplexer aus Aufgabe 4.1 zeichnen. Dieser besteht aus 2 AND und
		% einem OR Gatter. Auch hier ist es sinnvoll, aufgrund der Baum-Struktur von rechts anzufangen.
		%
		% Wir beginnen mit dem rechten OR-Gatter
		(0,0) node[or port, anchor=out] (OR) {}
		(OR.out) node[right] {Y}

		% Dann platzieren wir die beiden AND-Gatter Links davon, wobei das zweite einen negierten
		% Eingang bekommt
		(OR.in 1) -| ++ (0,0.5) node[and port, anchor=out] (AND1) {}
		(OR.in 2) -| ++ (0,-0.5) node[and port, anchor=out] (AND2) {}
		(AND2.bin 1) node[ocirc, left] {}

		% Dann beschriften wir die Eingänge A und B
		(AND1.in 1) -- ++ (-0.5,0) node[left] {A}
		(AND2.in 2) -- ++ (-0.5,0) node[left] {B}

		% Und hier beginnt das Problem: Wir müssen noch die "S" Leitung zeichnen. Diese hätten wir
		% gerne in der Mitte zwischen dem A und dem B Eingang. Außerdem muss ein Punkt dafür platziert
		% werden. Der Schönheit halber soll der Punkt ebenfalls mittig sein.
		%
		% Zuerst machen wir noch den einfachen Part: Wir verbinden die verbliebenen Eingänge
		% der beiden Gatter:
		(AND1.in 2) -- (AND2.in 1)
		% Jetzt hätten wir den Punkt gerne in der Mitte der beiden Eingänge.
		% Wir wissen: Die Mitte ist auf der Höhe von (OR.out). Der Punkt befindet sich also
		% an der X-Koordinate von (OR.out) und der Y-Koordinate von (AND1.in 2) oder (AND2.in 1).
		% Mit der nächsten Zeile lässt sich diese Koordinate bestimmen:
		(AND1.in 2 |- OR.out) node[circ] {}
		% Super. Von hier aus können wir noch die Beschriftung hinzufügen:
		-- ++ (-0.5,0) node[left] {S}
		% Der Ausdruck (A |- B) bestimmt eine Koordinate, die die X-Koordinate von A und die
		% Y-Koordinate von B hat.
		% Äquivalent dazu gibt es auch den Ausdruck (A -| B), der die Y-Koordinate von A und
		% die X-Koordinate von B bestimmt.
		%
		%Schnittpunkt-Koordinaten werden selten verwendet, sind aber ein wichtiges Werkzeug in Tikz.
	;
\end{circuitikz} \\
Bild 14: Ein Beispiel für Schnittpunkt-Koordinaten ist die Position des schwarzen Punktes

%%%%%%%%%%%%%%%%%%%%%%%%%%
% Kapitel 8: Multiplexer %
%%%%%%%%%%%%%%%%%%%%%%%%%%

% Multiplexer sind ebenfalls in er Circuitikz-Bibliothek enthalten. Im Vergleich zu allen
% anderen node-Elementen haben Multiplexer einen deutlich größeren Parameter-Umfang.
% In der Circuitikz-Anleitung werden sie in Kapitel 4.24 im Detail beschrieben.

\end{document}
